\chapter{Cơ sở lý thuyết}
\label{chap:theory}

\section{Tổng quan về mô hình ngôn ngữ lớn (LLM) và tối ưu ngữ cảnh}

Các mô hình ngôn ngữ lớn (Large Language Model - LLM) phát triển từ các kiến trúc mạng hồi quy (RNN) và LSTM truyền thống vốn bị giới hạn ở khả năng tính toán tuần tự và duy trì phụ thuộc tầm xa. Năm 2017, Vaswani và cộng sự đề xuất kiến trúc Transformer với cơ chế tự chú ý (self-attention), cho phép tính toán đồng thời mối liên hệ giữa mọi cặp phần tử trong chuỗi và khai thác triệt để khả năng song song hóa của phần cứng hiện đại. Trên nền tảng đó, các họ mô hình tự hồi quy (autoregressive) như GPT, LLaMA hay Qwen — được huấn luyện trên hàng nghìn tỷ token — đã thể hiện khả năng mô phỏng ngôn ngữ tự nhiên và suy luận ở mức độ đáng kể.

Về mặt kỹ thuật, LLM phân tách văn bản đầu vào thành các đơn vị cơ sở gọi là token thông qua thuật toán như Byte-Pair Encoding (BPE). Giới hạn cốt lõi của mô hình nằm ở \textbf{cửa sổ ngữ cảnh} (context window) — số lượng token tối đa được phép xử lý trong một lượt suy luận. Dù các kiến trúc thế hệ mới đã mở rộng giới hạn này lên hàng trăm nghìn token, nghiên cứu thực nghiệm cho thấy hiệu năng không tuyến tính với độ dài ngữ cảnh. Hiện tượng này được gọi là \textit{context rot}: khi số lượng token tăng lên, cơ chế tự chú ý phải duy trì $n^2$ quan hệ cặp trên toàn chuỗi, dẫn đến ngân sách chú ý (attention budget) bị phân tán và khả năng truy hồi thông tin suy giảm dần — ngay cả trước khi đạt ngưỡng giới hạn tuyệt đối. Điều này đặt ra yêu cầu tối ưu hoá chủ động toàn bộ nội dung đưa vào ngữ cảnh, thay vì chỉ tập trung vào cách diễn đạt câu lệnh.

Nhận thức này dẫn đến sự chuyển dịch từ \textbf{prompt engineering} sang \textbf{context engineering}. Prompt engineering tập trung vào việc thiết kế văn bản đầu vào nhằm tối ưu kết quả suy luận — chủ yếu thông qua system prompt, few-shot examples và các kỹ thuật như Chain-of-Thought. Đây là phương pháp phù hợp với các tác vụ đơn lượt (single-turn) hoặc phân loại văn bản đơn giản. Tuy nhiên, khi LLM vận hành trong vòng lặp tác nhân (agent loop) trải dài qua nhiều bước suy luận, lượng dữ liệu tích lũy — lịch sử hội thoại, kết quả gọi công cụ, thông tin truy vấn từ bộ nhớ ngoài — liên tục mở rộng và cạnh tranh nguồn tài nguyên chú ý. Anthropic định nghĩa context engineering là tập hợp các chiến lược tuyển chọn và duy trì bộ token tối ưu trong suốt quá trình suy luận, bao gồm system prompt, định nghĩa công cụ, dữ liệu truy xuất động và lịch sử tin nhắn. Nguyên tắc chỉ đạo là \textit{tìm tập hợp token nhỏ nhất có độ thông tin cao nhất} để tối đa hoá xác suất đạt kết quả mong muốn.

Trong kiến trúc tác nhân thông minh (AI Agent), ngữ cảnh được cấu thành từ hai thành phần cốt lõi:
\begin{itemize}
  \item \textbf{System Prompt (Chỉ dẫn hệ thống):} Thiết lập cấu hình nền tảng, xác định vai trò tác nhân, ràng buộc vận hành và heuristic hành vi xuyên suốt phiên làm việc. Một system prompt hiệu quả cần đạt \textit{đúng độ cao trừu tượng} (right altitude): đủ cụ thể để định hướng hành vi, nhưng đủ linh hoạt để mô hình suy luận tự chủ thay vì thực thi logic cứng nhắc.
  \item \textbf{User Prompt (Yêu cầu người dùng):} Mang nội dung tương tác trực tiếp của người dùng, đồng thời là điểm tích hợp dữ liệu ngữ cảnh động — bao gồm kết quả truy vấn bộ nhớ, đầu ra công cụ và thông tin truy xuất theo thời gian thực.
\end{itemize}

MiniClaw vận dụng các nguyên lý này thông qua cơ chế tổng hợp ngữ cảnh động tại mỗi lượt gọi mô hình: system prompt được xây dựng tự động dựa trên chế độ vận hành hiện tại, lịch sử hội thoại được kiểm soát theo quy tắc rút gọn xác định, và dữ liệu bộ nhớ được truy xuất có chọn lọc thay vì nạp toàn bộ vào ngữ cảnh — nhằm duy trì mật độ thông tin cao và hạn chế hiện tượng context rot.

\section{Kiến trúc Agent và vòng lặp tác nhân ReAct}

Thay vì phản hồi đơn lượt, kiến trúc tác nhân thông minh (Agentic Architecture) tổ chức quá trình suy luận theo một vòng lặp tác nhân (agent loop) — cho phép mô hình liên tục đánh giá trạng thái, lập kế hoạch và tương tác với môi trường bên ngoài cho đến khi hoàn thành mục tiêu.

Vòng lặp này được điều phối theo khung lập luận ReAct (Reasoning and Acting) \cite{yao2022react}, kết hợp chặt chẽ giữa suy luận logic và thực thi hành động qua ba giai đoạn lặp lại:
\begin{enumerate}
  \item \textbf{Thought (Lập luận):} LLM phân tích trạng thái hiện tại, xác định mục tiêu và hoạch định bước xử lý tiếp theo.
  \item \textbf{Action (Hành động):} Mô hình lựa chọn công cụ phù hợp từ tập định nghĩa sẵn — truy vấn cơ sở dữ liệu, ghi tệp, gọi API ngoại vi — kèm tham số tương ứng.
  \item \textbf{Observation (Quan sát):} Hệ thống thực thi công cụ, ghi nhận kết quả và đưa trở lại ngữ cảnh làm tiền đề cho lượt suy luận kế tiếp.
\end{enumerate}

Vòng lặp kết thúc khi LLM xác định đã đạt được mục tiêu và trả về phản hồi cuối cùng (\emph{Final Answer}). Hình~\ref{fig:agent-react-loop} mô tả dòng điều khiển tương ứng.

\begin{figure}[H]
  \centering
  \begin{tikzpicture}[
    node distance=1.8cm,
    block/.style={rectangle, draw, fill=blue!5, rounded corners, minimum width=3.8cm, minimum height=1cm, align=center, font=\small},
    decision/.style={diamond, draw, fill=orange!5, minimum width=3.2cm, minimum height=1cm, align=center, font=\small, inner sep=0pt},
    arrow/.style={thick,->,>=stealth}
  ]
    \node (user) [block] {Nhận yêu cầu\\từ người dùng};
    \node (thought) [block, below=0.5cm of user] {Quan sát và lập luận\\(Thought)};
    \node (decide) [decision, below=0.8cm of thought] {Quyết định\\hành động?};
    \node (action) [block, left=2.0cm of decide] {Thực thi công cụ\\(Action)};
    \node (obs) [block, above=1.6cm of action] {Quan sát kết quả\\(Observation)};
    \node (final) [block, below=0.8cm of decide] {Phản hồi\\người dùng};

    \draw [arrow] (user) -- (thought);
    \draw [arrow] (thought) -- (decide);
    \draw [arrow] (decide) -- node[anchor=south, font=\scriptsize] {Gọi công cụ} (action);
    \draw [arrow] (action) -- (obs);
    \draw [arrow] (obs) -- (thought.west);
    \draw [arrow] (decide) -- node[anchor=west, font=\scriptsize] {Hoàn thành} (final);
  \end{tikzpicture}
  \caption{\emph{Sơ đồ vòng lặp tác nhân ReAct kết hợp lập luận và hành động.}}
  \label{fig:agent-react-loop}
\end{figure}

Triển khai ReAct nguyên bản dựa trên sinh văn bản tự do để mã hóa lệnh gọi công cụ (ví dụ: \texttt{Action: write\_file[args]}), dẫn đến hai vấn đề thực tế: ambiguity trong phân tích cú pháp và hallucination cú pháp khi mô hình tự phát minh định dạng lệnh không tồn tại. Các hệ thống hiện đại khắc phục điều này bằng cơ chế \textbf{gọi công cụ có cấu trúc} (structured tool call): thay vì sinh văn bản tự do, LLM tiếp nhận mô tả công cụ chuẩn hóa theo JSON Schema và xuất trực tiếp đối tượng JSON chứa định danh hàm cùng tham số. Một lớp điều phối công cụ (Tool Harness) chịu trách nhiệm xác thực kiểu, định tuyến thực thi và xử lý ngoại lệ — tách biệt hoàn toàn nghiệp vụ khỏi lớp suy luận của mô hình.

MiniClaw hiện thực hóa kiến trúc này với Tool Harness được xây dựng trên TypeScript, tận dụng hệ thống kiểu tĩnh để đảm bảo tính đúng đắn của tham số ngay tại compile-time trước khi thực thi.

\section{Quản lý bộ nhớ tác nhân (Agent Memory)}
Bộ nhớ tác nhân đóng vai trò duy trì tính nhất quán hành vi qua các phiên làm việc kéo dài, đồng thời hỗ trợ cá nhân hóa trải nghiệm người dùng. Hệ thống phân tầng kiến trúc bộ nhớ này thành hai lớp chính:

\subsection{Bộ nhớ ngắn hạn (short-term memory)}
Bộ nhớ ngắn hạn duy trì ngữ cảnh tương tác tức thời trong phiên làm việc hiện tại. Hệ thống lưu trữ toàn bộ chuỗi thông điệp (truy vấn từ người dùng, phản hồi của tác nhân, các yêu cầu gọi công cụ và kết quả thực thi tương ứng) dưới dạng một danh sách tuyến tính có thứ tự. Danh sách này tích hợp trực tiếp vào cửa sổ ngữ cảnh của LLM tại mỗi lượt gọi API kế tiếp. Khi tần suất hội thoại tăng, dung lượng bộ nhớ ngắn hạn gia tăng sẽ gây áp lực lớn lên giới hạn của cửa sổ ngữ cảnh.

\subsection{Bộ nhớ dài hạn (long-term memory)}
Bộ nhớ dài hạn lưu trữ các sự kiện, thông tin cấu hình và thói quen mang tính bền vững của người dùng qua các phiên làm việc. Hệ thống không thể đưa toàn bộ bộ nhớ dài hạn vào prompt vì giới hạn token và chi phí tính toán. Thay vào đó, hệ thống phân tách thông tin dài hạn thành các phân đoạn độc lập và lưu trữ trong cơ sở dữ liệu vector (vector database) ngoại vi \cite{memverge2026memmachine}.

Quy trình truy xuất dữ liệu dài hạn chuyển đổi các phân đoạn ngôn ngữ thành vector số thực đa chiều nhờ mô hình nhúng văn bản (Text Embedding). Khi tiếp nhận yêu cầu mới, hệ thống chuyển truy vấn thành vector đại diện $u$, sau đó tìm kiếm tương đồng ngữ nghĩa trên không gian các vector ký ức $v$ bằng độ tương đồng Cosine (Cosine Similarity):
\begin{equation}
  \text{CosineSimilarity}(u, v) = \frac{u \cdot v}{\|u\| \|v\|} = \frac{\sum_{i=1}^{n} u_i v_i}{\sqrt{\sum_{i=1}^{n} u_i^2} \sqrt{\sum_{i=1}^{n} v_i^2}}
\end{equation}
Hệ thống tự động trích xuất các phân đoạn ký ức có chỉ số tương đồng cao nhất vượt qua ngưỡng thiết lập để bổ sung vào prompt làm ngữ cảnh tham chiếu. Cơ chế này giúp tác nhân tái hiện thông tin lịch sử một cách có chọn lọc \cite{yu2026agentic}. Hình \ref{fig:memory-hierarchy} mô tả chi tiết dòng chảy dữ liệu của quy trình này.

\begin{figure}[H]
  \centering
  \begin{tikzpicture}[
    node distance=1.5cm,
    block/.style={rectangle, draw, fill=blue!5, rounded corners, minimum width=3.6cm, minimum height=1.0cm, align=center, font=\small},
    database/.style={
      cylinder, draw, fill=gray!10, 
      shape border rotate=90, 
      aspect=0.4, 
      minimum width=2.4cm, 
      minimum height=1.4cm, 
      align=center, font=\small, text width=2.2cm
    },
    arrow/.style={thick,->,>=stealth}
]
    \node (query) [block] {Yêu cầu\\hội thoại mới};
    \node (embed) [block, below=0.6cm of query] {Mô hình nhúng và\\tìm kiếm vector};
    \node (vector) [database, right=1.8cm of embed] {Cơ sở dữ liệu\\vector};
    
    \node (mixer) [block, below=1.0cm of embed] {Bộ trộn ngữ cảnh\\(Context mixer)};
    \node (short) [block, left=1.5cm of mixer] {Lịch sử hội thoại\\ngắn hạn};
    \node (system) [block, right=1.5cm of mixer] {Chỉ dẫn hệ thống\\(System prompt)};
    
    \node (llm) [block, below=0.8cm of mixer] {Cửa sổ ngữ cảnh\\LLM};

    \draw [arrow] (query) -- (embed);
    \draw [arrow] (embed) -- node[anchor=south, font=\scriptsize] {Tìm kiếm Cosine} (vector);
    
    \draw [arrow] (vector.south) -| (mixer.north east);
    \draw [arrow] (short) -- (mixer);
    \draw [arrow] (system) -- (mixer);
    
    \draw [arrow] (mixer) -- (llm);
  \end{tikzpicture}
  \caption[\emph{Sơ đồ cấu trúc phân tầng bộ nhớ.}]{\emph{Sơ đồ cấu trúc phân tầng bộ nhớ ngắn hạn và dài hạn của tác nhân.}}
  \label{fig:memory-hierarchy}
\end{figure}

Sơ đồ Hình \ref{fig:memory-hierarchy} mô tả chi tiết kiến trúc phân tầng bộ nhớ. Sự phân tách này giải quyết bài toán cân đối giữa độ chính xác ngữ cảnh ngắn hạn và chi phí xử lý dài hạn. Bộ nhớ ngắn hạn lưu giữ chi tiết cuộc trò chuyện hiện tại nhưng bị giới hạn bởi dung lượng cửa sổ ngữ cảnh của mô hình. Trong khi đó, bộ nhớ dài hạn chỉ truy xuất những ký ức liên quan qua tìm kiếm vector tương đồng Cosine, giúp tối ưu lượng token tiêu thụ mà vẫn đảm bảo tính cá nhân hóa.

MiniClaw hiện thực hóa kiến trúc bộ nhớ phân tầng này thông qua việc sử dụng tệp lưu vết checkpoint cục bộ cho bộ nhớ ngắn hạn. Đồng thời, hệ thống triển khai mô hình nhúng cục bộ kết hợp công cụ tìm kiếm tương đồng Cosine trên cơ sở dữ liệu vector để truy xuất thông tin dài hạn phù hợp với ngữ cảnh hiện tại.

\section{Tối ưu ngữ cảnh và cơ chế nén ngữ cảnh tự động}
Các hệ thống tác nhân thông minh vận hành theo mô hình kiểm soát trạng thái (stateful) đòi hỏi kiểm soát dung lượng ngữ cảnh chặt chẽ. Do giới hạn cố định của cửa sổ ngữ cảnh LLM, chi phí tính toán và độ trễ phản hồi (latency) tăng tỷ lệ thuận với khối lượng dữ liệu đầu vào. Do đó, hệ thống cần áp dụng các giải pháp tối ưu ngữ cảnh nhằm quản lý hiệu quả ngân sách token (token budgeting).

Đồ án triển khai cơ chế nén ngữ cảnh (context compaction) tự động nhằm giải phóng tài nguyên bộ nhớ ngắn hạn. Với cấu hình thử nghiệm, hệ thống thiết lập ngưỡng kích hoạt nén tại mức \textbf{50.000 token}. Khi tổng độ dài lịch sử thông điệp vượt quá ngưỡng này, hệ thống tự động khởi chạy tiến trình nén ngữ cảnh ngầm theo các giai đoạn sau:
\begin{enumerate}
  \item \textbf{Phân tách lịch sử:} Hệ thống giữ lại một số lượng thông điệp gần nhất (như 10 lượt tương tác cuối) để duy trì mạch hội thoại trực tiếp.
  \item \textbf{Tóm tắt ngữ cảnh:} Hệ thống chuyển tiếp các thông điệp cũ hơn đến phân hệ LLM tóm tắt chuyên biệt để trích xuất thông tin cốt lõi.
  \item \textbf{Cập nhật hồ sơ dữ liệu:} Phân hệ LLM phân tích và cập nhật hai luồng dữ liệu cấu trúc:
    \begin{itemize}
      \item Bản ghi tóm tắt súc tích các sự kiện chính hoặc tác vụ đã xử lý thành công.
      \item Thông tin cập nhật bổ sung vào hồ sơ người dùng (User Profile), như các thiết lập lịch trình mới.
    \end{itemize}
  \item \textbf{Giải phóng bộ lưu vết (Checkpointer):} Hệ thống đồng bộ kết quả tóm tắt mới vào bộ lưu trữ trạng thái lâu dài, đồng thời giải phóng lịch sử thông điệp thô khỏi bộ nhớ checkpointer. Giai đoạn này tối ưu dung lượng bộ nhớ ngắn hạn nhưng vẫn đảm bảo tính toàn vẹn của thông tin cốt lõi.
\end{enumerate}
Quy trình điều phối này được trực quan hóa qua sơ đồ Hình \ref{fig:context-compaction}.

\begin{figure}[H]
  \centering
  \begin{tikzpicture}[
    node distance=1.4cm,
    block/.style={rectangle, draw, fill=blue!5, rounded corners, minimum width=3.8cm, minimum height=1cm, align=center, font=\small},
    decision/.style={diamond, draw, fill=orange!5, minimum width=3.2cm, minimum height=1cm, align=center, font=\small, inner sep=0pt, text width=2.4cm},
    arrow/.style={thick,->,>=stealth}
  ]
    \node (inbound) [block] {Tin nhắn\\hội thoại mới};
    \node (check) [decision, below=0.5cm of inbound] {Độ dài lịch sử\\> 50.000 token?};
    \node (normal) [block, right=1.8cm of check] {Lưu trữ và\\xử lý bình thường};
    \node (extract) [block, below=0.8cm of check] {Trích xuất\\tin nhắn cũ};
    \node (llm) [block, below=0.5cm of extract] {Mô hình tóm tắt\\và cập nhật hồ sơ};
    \node (save) [block, below=0.5cm of llm] {Lưu hồ sơ\\và bản tóm tắt mới};
    \node (prune) [block, below=0.5cm of save] {Giải phóng lịch sử\\tin nhắn thô};
    \node (done) [block, below=0.6cm of prune] {Hoàn thành\\điều phối ngữ cảnh};

    \draw [arrow] (inbound) -- (check);
    \draw [arrow] (check) -- node[anchor=south, font=\scriptsize] {Không} (normal);
    \draw [arrow] (check) -- node[anchor=west, font=\scriptsize] {Có} (extract);
    \draw [arrow] (extract) -- (llm);
    \draw [arrow] (llm) -- (save);
    \draw [arrow] (save) -- (prune);
    \draw [arrow] (prune) -- (done);
    
    \draw [arrow] (normal.south) |- ([yshift=0.3cm]done.north) -- (done.north);
  \end{tikzpicture}
  \caption[\emph{Quy trình điều phối và nén ngữ cảnh.}]{\emph{Quy trình điều phối và nén ngữ cảnh tự động tối ưu hóa ngân sách token.}}
  \label{fig:context-compaction}
\end{figure}

Sơ đồ Hình \ref{fig:context-compaction} thể hiện rõ ràng cơ chế phân nhánh để bảo vệ cửa sổ ngữ cảnh. Thiết kế này giải quyết trade-off giữa độ chi tiết của lịch sử hội thoại và tốc độ phản hồi của hệ thống. Việc nén ngầm các thông điệp cũ giúp duy trì tốc độ xử lý nhanh, đồng thời tránh hiện tượng mất mát thông tin quan trọng nhờ việc tóm tắt thông tin và cập nhật vào hồ sơ người dùng một cách có cấu trúc.

MiniClaw hiện thực hóa nguyên lý nén ngữ cảnh tự động này bằng cách liên tục giám sát số lượng token tích lũy qua mỗi lượt hội thoại. Khi bộ đếm đạt ngưỡng giới hạn, hệ thống kích hoạt tiến trình nén ngầm. Tiến trình này cập nhật dữ liệu cấu trúc vào hồ sơ người dùng mà không gây gián đoạn luồng tương tác chính.

\section{Trích xuất quy trình làm việc có kiểm soát (HITL)}
Khi giải quyết các nhiệm vụ phức tạp, tác nhân thường gọi nhiều công cụ liên tục để tìm kiếm thông tin và khám phá ngữ cảnh (skill/context discovery). Việc lặp lại chuỗi công cụ dài này cho các yêu cầu tương tự gây lãng phí token, tăng độ trễ và tích lũy sai số. Để tối ưu hiệu năng, đồ án đề xuất cơ chế trích xuất quy trình tái sử dụng (Reusable Workflow Extraction). Cơ chế này tự động nhận diện và đóng gói chuỗi hành động thành một kỹ năng (skill) mới. Gặp bài toán tương tự, tác nhân sẽ gọi trực tiếp quy trình đã lưu thay vì chạy lại các bước tìm kiếm. Nhằm đảm bảo tính chính xác và an toàn cho hệ thống, đồ án tích hợp cơ chế kiểm soát bởi người dùng (Human-in-the-Loop - HITL) để phê duyệt và hiệu chỉnh logic trước khi lưu trữ.

Hệ thống phát hiện các chuỗi hành động lặp lại dựa trên phân tích lịch sử tương tác (ví dụ: truy vấn dữ liệu định kỳ hay lọc thông tin). Thay vì yêu cầu người dùng thiết lập thủ công, tác nhân đề xuất đóng gói trình tự đó thành một kỹ năng có khả năng tái sử dụng.

Quy trình tương tác HITL này vận hành qua ba bước chính:
\begin{enumerate}
  \item \textbf{Đề xuất quy trình:} Hệ thống tạo mã nguồn hoặc chỉ dẫn cho quy trình mới dưới dạng tệp tài liệu chuẩn hóa \texttt{SKILL.md} để lưu tạm thời.
  \item \textbf{Yêu cầu phê duyệt:} Tác nhân mở một luồng phụ thông báo cấu trúc quy trình đề xuất cho người dùng và tiếp nhận phản hồi qua Telegram Bot API.
  \item \textbf{Nạp động quy trình (Dynamic Loading):} 
    \begin{itemize}
      \item Khi người dùng chấp thuận (\emph{Approve}), hệ thống di chuyển tệp cấu trúc sang thư mục lưu trữ chính thức \texttt{workflows}. Sau đó, hệ thống tự động quét và nạp động (dynamic loading) cấu hình mới vào danh mục công cụ khả dụng của tác nhân mà không cần khởi động lại máy chủ.
      \item Khi người dùng từ chối (\emph{Reject}), hệ thống hủy bỏ bản nháp và khôi phục trạng thái vận hành ban đầu.
    \end{itemize}
\end{enumerate}
Sơ đồ Hình \ref{fig:hitl-feedback-loop} thể hiện toàn bộ chu trình phản hồi này.

\begin{figure}[H]
  \centering
  \begin{tikzpicture}[
    node distance=1.4cm,
    block/.style={rectangle, draw, fill=blue!5, rounded corners, minimum width=3.8cm, minimum height=1cm, align=center, font=\small},
    decision/.style={diamond, draw, fill=orange!5, minimum width=3.2cm, minimum height=1cm, align=center, font=\small, inner sep=0pt, text width=2.4cm},
    arrow/.style={thick,->,>=stealth}
  ]
    \node (detect) [block] {Phát hiện mẫu\\hành vi lặp lại};
    \node (write) [block, below=0.5cm of detect] {Trích xuất quy trình\\thành tệp \texttt{SKILL.md}};
    \node (telegram) [block, below=0.5cm of write] {Gửi yêu cầu\\phê duyệt Telegram};
    \node (approve) [decision, below=0.6cm of telegram] {Người dùng\\chấp thuận?};
    
    \node (reject) [block, left=1.8cm of approve] {Hủy bỏ bản nháp\\đề xuất};
    
    \node (save) [block, below=1.0cm of approve] {Lưu tệp quy trình\\vào thư mục lưu trữ};
    \node (load) [block, below=0.5cm of save] {Nạp động quy trình\\và kích hoạt công cụ};
    
    \node (restore) [block, below=0.8cm of load] {Khôi phục luồng\\hoạt động chính};

    \draw [arrow] (detect) -- (write);
    \draw [arrow] (write) -- (telegram);
    \draw [arrow] (telegram) -- (approve);
    
    \draw [arrow] (approve) -- node[anchor=south, font=\scriptsize] {Từ chối} (reject);
    \draw [arrow] (approve) -- node[anchor=west, font=\scriptsize] {Đồng ý} (save);
    
    \draw [arrow] (save) -- (load);
    \draw [arrow] (load) -- (restore);
    
    \draw [arrow] (reject.south) |- ([yshift=0.4cm]restore.north) -- (restore.north);
  \end{tikzpicture}
  \caption[\emph{Sơ đồ vòng lặp phản hồi HITL.}]{\emph{Sơ đồ vòng lặp phản hồi kiểm soát bởi người dùng (HITL) phê duyệt quy trình làm việc mới.}}
  \label{fig:hitl-feedback-loop}
\end{figure}

Sơ đồ Hình \ref{fig:hitl-feedback-loop} minh họa rõ ràng cơ chế kiểm soát có phản hồi từ người dùng. Thiết kế này giải quyết trade-off giữa tính tự động hóa và độ an toàn của hệ thống. Dù việc nạp động trực tiếp giúp nâng cao độ linh hoạt, hệ thống vẫn bắt buộc đi qua bước phê duyệt trung gian nhằm loại bỏ hoàn toàn rủi ro thực thi mã nguồn lỗi hoặc hành động ngoài ý muốn do ảo giác mô hình (hallucination).

MiniClaw ứng dụng cơ chế trích xuất quy trình làm việc có kiểm soát (HITL) này để tự động cập nhật mã nguồn kỹ năng một cách an toàn. Việc tách biệt giai đoạn đề xuất lưu vết tạm thời trong tệp \texttt{SKILL.md} và nạp động từ thư mục \texttt{workflows} sau khi người dùng phê duyệt qua Telegram giúp bảo vệ tính toàn vẹn của hệ thống.

\section{Khung công cụ tác nhân (Agent Tool Harness)}
Khung công cụ tác nhân (Agent Tool Harness) đóng vai trò là giao tiếp trung gian, kết nối phân hệ suy luận LLM với môi trường hệ điều hành và các dịch vụ mạng ngoại vi. Thành phần này chuyển hóa đầu ra suy luận trừu tượng của mô hình thành thao tác dữ liệu hoặc tương tác hệ thống cụ thể.

\subsection{Khái niệm khung công cụ và vai trò trong vòng lặp tác nhân}
Hệ thống đặc tả mỗi công cụ (Tool) qua hai thành phần bắt buộc:
\begin{itemize}
  \item \textbf{JSON Schema:} Bản mô tả chi tiết gồm định danh, chức năng nghiệp vụ bằng ngôn ngữ tự nhiên để hỗ trợ LLM định tuyến chính xác, cùng cấu trúc kiểu dữ liệu của tham số đầu vào.
  \item \textbf{Hàm thực thi (Implementation function):} Hàm xử lý mã nguồn TypeScript tiếp nhận các tham số đã qua xác thực từ mô hình để tiến hành xử lý logic nghiệp vụ.
\end{itemize}
Vòng lặp tác nhân tích hợp khung công cụ (Tool Harness) vận hành theo cơ chế thống nhất. LLM tiếp nhận toàn bộ ngữ cảnh phiên cùng danh sách cấu trúc các công cụ khả dụng. Khi mô hình phát yêu cầu gọi công cụ (Tool Call), bộ điều phối bắt giữ thông điệp, ánh xạ định danh với hàm thực thi tương ứng, xử lý và trả kết quả về không gian ngữ cảnh của mô hình.

\subsection{LangGraph và trừu tượng hóa đồ thị trạng thái}
Thư viện LangGraph cung cấp mô hình trừu tượng hóa vòng lặp tác nhân dưới dạng đồ thị trạng thái (State Graph). Trong kiến trúc này, các nút (Nodes) biểu diễn các bước xử lý dữ liệu và chia sẻ trạng thái chung của hệ thống (gồm lịch sử thông điệp và biến trạng thái nội bộ), còn các cạnh (Edges) chịu trách nhiệm điều phối dòng chảy:
\begin{itemize}
  \item \textbf{Nút Agent (Agent Node):} Chứa mô hình ngôn ngữ lớn, tiếp nhận trạng thái hiện tại của đồ thị để đưa ra phản hồi trực tiếp hoặc phát yêu cầu gọi công cụ.
  \item \textbf{Nút công cụ (Tools Node):} Tiếp nhận các yêu cầu gọi hàm, thực thi song song hoặc tuần tự các tiến trình tương ứng trong khung công cụ (Tool Harness), thu thập kết quả và cập nhật vào dòng trạng thái đồ thị.
  \item \textbf{Cạnh điều kiện (Conditional Edge):} Kiểm tra xem thông điệp cuối cùng trong trạng thái đồ thị có chứa yêu cầu gọi công cụ hay không. Nếu có, hệ thống định tuyến luồng xử lý sang nút công cụ (Tools Node); ngược lại, luồng xử lý dẫn hướng về điểm kết thúc (\texttt{\_\_end\_\_}).
\end{itemize}
Quy trình này được thể hiện rõ nét qua sơ đồ Hình \ref{fig:langgraph-state-graph}.

\begin{figure}[H]
  \centering
  \begin{tikzpicture}[
    node distance=1.5cm,
    block/.style={rectangle, draw, fill=blue!5, rounded corners, minimum width=3.5cm, minimum height=1cm, align=center, font=\small},
    decision/.style={diamond, draw, fill=orange!5, minimum width=3.2cm, minimum height=1.0cm, align=center, font=\small, inner sep=0pt, text width=2.2cm},
    point/.style={circle, draw, fill=red!10, minimum size=1cm, font=\small},
    arrow/.style={thick,->,>=stealth}
  ]
    \node (start) [point] {\texttt{\_\_start\_\_}};
    \node (agent) [block, below=0.8cm of start] {Nút Agent\\(Agent node)};
    \node (cond) [decision, below=0.8cm of agent] {Kiểm tra\\gọi công cụ?};
    \node (tools) [block, left=2.0cm of cond] {Nút công cụ\\(Tools node)};
    \node (end) [point, below=0.8cm of cond] {\texttt{\_\_end\_\_}};

    \draw [arrow] (start) -- (agent);
    \draw [arrow] (agent) -- (cond);
    
    \draw [arrow] (cond) -- node[anchor=south, font=\scriptsize] {Có yêu cầu} (tools);
    \draw [arrow] (cond) -- node[anchor=west, font=\scriptsize] {Không} (end);
    
    \draw [arrow] (tools.north) |- ([yshift=0.4cm]agent.north) -- (agent.north);
  \end{tikzpicture}
  \caption[\emph{Sơ đồ vòng lặp tác nhân LangGraph state graph.}]{\emph{Sơ đồ trừu tượng hóa vòng lặp tác nhân dựa trên đồ thị trạng thái LangGraph.}}
  \label{fig:langgraph-state-graph}
\end{figure}

Sơ đồ Hình \ref{fig:langgraph-state-graph} chỉ ra cách đồ thị trạng thái LangGraph quản lý dòng dữ liệu và kiểm soát vòng lặp. Mô hình này mang lại khả năng phân tách rõ ràng giữa phần suy luận của mô hình (Agent Node) và phần thực thi mã lệnh (Tools Node). Mặc dù cấu trúc này tăng nhẹ độ phức tạp khi lập trình đồ thị, nó mang lại tính ổn định cao và khả năng bảo trì tốt hơn cho hệ thống.

\subsection{Phân loại công cụ và ranh giới bảo mật}
Hệ thống phân lớp công cụ chặt chẽ dựa trên mức độ tác động của chúng đối với tài nguyên máy chủ để xác lập ranh giới kiểm soát an toàn:
\begin{enumerate}
  \item \textbf{Chỉ đọc (Read-only):} Các tác vụ truy vấn dữ liệu không thay đổi trạng thái hệ thống, như đọc tệp tài liệu cấu hình, tìm kiếm ngữ nghĩa hoặc truy vấn cơ sở dữ liệu vector.
  \item \textbf{Ghi dữ liệu (Write):} Các tác vụ thay đổi trạng thái lưu trữ nội bộ, bao gồm tạo mới lịch trình hoặc cập nhật danh mục ghi nhớ cá nhân.
  \item \textbf{Có tác động phụ (Side-effect / System actions):} Các tác vụ tương tác trực tiếp với hệ điều hành máy chủ và Internet, như gửi email, tải xuống tài nguyên hoặc thực thi các câu lệnh CLI.
\end{enumerate}

Để đảm bảo an toàn thực thi và ngăn chặn rủi ro tiêm nhiễm chỉ thị (prompt injection) dẫn đến kích hoạt mã độc, khung công cụ thiết lập cơ chế cô lập nghiêm ngặt. Hệ thống vận hành các tác vụ nguy cơ cao trong môi trường hộp cát (sandbox) hạn chế đặc quyền, đồng thời áp dụng các tầng lọc và xác thực tham số đầu vào trước khi chuyển đến tầng thực thi lõi.

Hệ thống MiniClaw sử dụng lớp \texttt{FileCheckpointSaver} để ghi nhận trạng thái đồ thị. Dựa vào kết quả định tuyến từ đồ thị LangGraph, hệ thống sẽ gọi các công cụ tương ứng.

\section{Các công nghệ và thư viện phát triển chính}
Để xây dựng hệ thống tác nhân tối ưu tài nguyên trên hạ tầng sẵn có và đảm bảo khả năng mở rộng, đồ án lựa chọn và tích hợp các công nghệ cốt lõi sau:

\subsection{Node.js và TypeScript}
Hệ thống sử dụng nền tảng backend Node.js nhờ cơ chế xử lý bất đồng bộ và hướng sự kiện (event-driven, non-blocking I/O). Đặc tính này giúp tối ưu hóa việc truyền nhận dữ liệu và phản hồi theo thời gian thực của chatbot.

Ngôn ngữ TypeScript hỗ trợ kiểm tra kiểu tĩnh (static typing) nghiêm ngặt để quản lý mã nguồn chặt chẽ. Hệ thống phát hiện và ngăn chặn sớm các lỗi cấu trúc dữ liệu ngay trong giai đoạn biên dịch (compile-time), tránh lỗi vận hành (runtime) và tăng tính mô-đun hóa để dễ dàng bảo trì.

\subsection{Grammy — Thư viện Telegram Bot}
Đồ án lựa chọn Grammy làm thư viện chính để giao tiếp với Telegram Bot API. Grammy đạt hiệu năng xử lý bất đồng bộ cao trong môi trường Node.js và hỗ trợ tốt cho TypeScript. Với kiến trúc phân tầng mã trung gian (middleware) mạnh mẽ, Grammy giúp hệ thống triển khai các bộ lọc thông điệp, quản lý phiên trò chuyện (session) và định dạng giao diện tương tác linh hoạt.

\subsection{LangGraph — Khung điều phối đồ thị trạng thái}
Để quản lý logic suy luận đa bước của tác nhân, đồ án áp dụng LangGraph làm khung điều phối (orchestration engine) theo mô hình đồ thị trạng thái (State Graph). Giải pháp này giúp kiểm soát chi tiết tiến trình lập luận (reasoning), điều hướng linh hoạt giữa các trạng thái và xử lý các yêu cầu gọi công cụ song song.

Đặc biệt, tính năng lưu vết điểm kiểm soát (checkpointer) của LangGraph cung cấp cho hệ thống khả năng tự phục hồi trạng thái (fault tolerance) khi gặp sự cố gián đoạn hoặc ngoại lệ.

\subsection{Ollama và các nhà cung cấp LLM ngoại vi}
Tầng mô hình ngôn ngữ lớn (LLM) kết hợp linh hoạt giữa giải pháp cục bộ và đám mây nhằm cân đối hiệu năng, chi phí và tính bảo mật:
\begin{itemize}
    \item \textbf{Ollama:} Hỗ trợ chạy cục bộ (self-hosted) các mô hình mã nguồn mở (như LLaMA hoặc Qwen) trên phần cứng giới hạn. Ollama giúp bảo vệ quyền riêng tư của dữ liệu và đảm bảo hệ thống hoạt động độc lập không phụ thuộc Internet.
    \item \textbf{API đám mây thương mại:} Đóng vai trò là cơ chế dự phòng (fallback). Khi gặp tác vụ phức tạp vượt quá năng lực suy luận của mô hình cục bộ, hệ thống tự động chuyển hướng gọi API đám mây nhằm bảo đảm chất lượng phản hồi.
\end{itemize}
Sự phân cấp công nghệ này được trực quan hóa qua sơ đồ Hình \ref{fig:technology-stack}.

\begin{figure}[H]
  \centering
  \begin{tikzpicture}[
    node distance=1.1cm,
    layer/.style={rectangle, draw, fill=blue!5, rounded corners=2pt, minimum width=\linewidth, text width=0.96\linewidth, minimum height=0.8cm, align=center, font=\small},
    title/.style={font=\bfseries\small}
  ]
    % Định nghĩa các tầng công nghệ trong hệ thống
    \node (ui) [layer, fill=teal!10] {\textbf{Tầng giao tiếp:} Telegram Bot API và thư viện Grammy};
    \node (logic) [layer, below=0.4cm of ui, fill=blue!5] {\textbf{Tầng điều phối:} Đồ thị trạng thái LangGraph};
    \node (engine) [layer, below=0.4cm of logic, fill=orange!5] {\textbf{Tầng suy luận và công cụ:} Khung công cụ, Ollama (LLaMA, Qwen) và API đám mây};
    \node (base) [layer, below=0.4cm of engine, fill=gray!10] {\textbf{Nền tảng thực thi:} Môi trường Node.js và TypeScript};

    % Vẽ các mũi tên tương tác hai chiều giữa các tầng
    \draw [thick,<->,>=stealth] (ui) -- (logic);
    \draw [thick,<->,>=stealth] (logic) -- (engine);
    \draw [thick,<->,>=stealth] (engine) -- (base);
  \end{tikzpicture}
  \caption[\emph{Sơ đồ phân tầng tích hợp công nghệ MiniClaw.}]{\emph{Sơ đồ phân tầng cấu trúc tích hợp công nghệ trong hệ thống tác nhân.}}
  \label{fig:technology-stack}
\end{figure}

Sơ đồ Hình \ref{fig:technology-stack} phân tầng rõ ràng kiến trúc hệ thống, từ giao diện người dùng đến tầng thực thi mã nguồn. Sự kết hợp này mang lại khả năng phân tách trách nhiệm (separation of concerns), giúp bảo trì và nâng cấp từng tầng độc lập. Tuy nhiên, việc tích hợp nhiều công nghệ đòi hỏi cơ chế quản lý trạng thái đồng bộ chặt chẽ nhằm tránh xung đột dữ liệu giữa tầng điều phối đồ thị và tầng thực thi công cụ.

Tóm lại, Chương II đã hệ thống hóa toàn bộ cơ sở lý thuyết cần thiết bao gồm kiến trúc Transformer, vòng lặp tác nhân ReAct, các tầng bộ nhớ, cơ chế nén ngữ cảnh và khung công cụ. Những nguyên lý này đóng vai trò nền tảng để giải quyết các bài toán về giới hạn token, tính riêng tư và khả năng tự động hóa của hệ thống. Đây là tiền đề trực tiếp để đồ án tiến hành thiết kế kiến trúc chi tiết và xây dựng các chức năng thực tiễn của trợ lý ảo MiniClaw trong Chương III.